%Roman numeral page numbers for the front material.
\phantomsection
\addcontentsline{toc}{chapter}{Editor's Preface}
\fancyhead[C]{{\LARGE Editor's Preface}}

\begin{center}
	`Septies in die laudem dixi tibi : super judicia justicie tue.'
\end{center}

%\lett{T}{he} laudable practice of praying the Daily Office has increasingly been received and recovered by the People of God in these our times, and with a great fervour. For most, the Book of Common is a sufficient resource and pattern of life. However, we have received increasing requests for all the Hours of the Church to be published with the full propers throughout the Year. In love for Christ and consideration of our patrons and fellow workers, we are happy to publish this Breviary according to the Anglican Rite of the Catholic Church. We are confident that it will prove to be a helpful resource for all who drink from the rich fount of Christ.\\

\begin{otherlanguage}{ecclesiasticallatin}
\par\noindent
\lett{L}{audabilis} consuetudo illa cotidiani Officii persolvendi his nostris temporibus a Populo Dei in dies magis recepta est atque recuperata, idque fervore non modico. Plerisque quidem Liber Precum Publicarum et sufficiens adiumentum et forma vivendi existit. Verumtamen crebrescentes ad nos peticiones pervenerunt, ut omnes Ecclesie Hore cum integris per tocius Anni circulum propriis ederentur. Christi igitur amore, necnon benefactorum ac cooperatorum nostrorum intuitu, hoc Breviarium secundum Ritum Anglicanum Ecclesie Catholice gratanter in lucem emittimus. Confidimus autem quod omnibus qui de fonte Christi uberrimo bibunt utile subsidium sit futurum.\\

\end{otherlanguage}


\rightline{On the Feast of the Conception of the Blessed Virgin Mary}

\vspace{1\baselineskip}

\rightline{\emph{Pax Christi,}}

\rightline{Mr. J. P. Augustine Watson, M.St.}

\rightline{\emph{Chief Editor, Apologia Anglicana}}

\rightline{\emph{Ecclesiarch, St. Botolph Orthodox Chapel}}